% Abstract — two concise paragraphs, standalone page
\section*{Abstract}
\addcontentsline{toc}{section}{Abstract}

Academic research workflows remain heavily fragmented, forcing students and scholars to juggle disconnected tools for paper storage, PDF reading, note-taking, reference management, and manuscript writing. A requirements survey among academic researchers confirmed substantial productivity loss due to context switching, low satisfaction with legacy reference managers, and universal demand for integrated artificial intelligence assistance. This report presents \textbf{ScholarFlow}, a full-stack, AI-powered research paper collaboration hub designed to unify the end-to-end research lifecycle into a single collaborative platform: cloud document management with automated metadata extraction, semantic search powered by vector embeddings, inline PDF annotations, collaborative TipTap academic writing, and multi-tenant research workspaces.

The platform is engineered as a decoupled TypeScript system comprising a Next.js 16 frontend, an Express.js REST and WebSocket backend, a PostgreSQL database enhanced with pgvector, and an AWS S3 storage infrastructure with presigned access. Integrated multi-model AI orchestration powers contextual summarization, structured Q\&A insights, and multi-paper synthesis under strict workspace access controls, while production-grade billing is managed via idempotent Stripe webhooks. ScholarFlow demonstrates the practical application of software engineering principles---modular design, automated testing (21 test suites / 59 tests), strict type safety, and real-time synchronization---to deliver an accessible, scalable, and secure platform for modern academic collaboration.
